% !TeX program = xelatex
\documentclass[11pt,a4paper]{article}
\usepackage[margin=23mm]{geometry}
\usepackage{fontspec}
\setmainfont{Latin Modern Roman}
\usepackage{amsmath,amssymb,mathtools}
\usepackage{xcolor,enumitem,booktabs,longtable,array,xurl,hyperref}
\definecolor{muted}{HTML}{53636C}
\hypersetup{colorlinks=true,urlcolor=blue,linkcolor=blue}
\setlength{\parindent}{0pt}
\setlength{\parskip}{5pt}
\newcommand{\R}{\mathbb R}
\newcommand{\E}{\mathbb E}
\newcommand{\N}{\mathbb N}
\newcommand{\ind}{\mathbf 1}
\newcommand{\Var}{\operatorname{Var}}
\newcommand{\Cov}{\operatorname{Cov}}
\newcommand{\supp}{\operatorname{supp}}
\DeclareMathOperator*{\argmin}{arg\,min}
\DeclareMathOperator*{\argmax}{arg\,max}
\newcommand{\entrymeta}[3]{{\sffamily\small\color{muted}\textbf{\texttt{#1}}\quad|\quad #2\par #3\par}}
\newcommand{\Problemref}[1]{\texttt{#1}}
\newcommand{\JELref}[2]{\texttt{#2} (JEL \texttt{#1})}
\begin{document}
\section*{Climate uncertainty in dynamic spatial welfare and policy}
\textbf{Problem ID:} \texttt{Q54-88}\quad \textbf{JEL:} \texttt{Q54}\par
% BEGIN ECONOMICS STATEMENT
\label{problem:OP-0615}


\label{audited:ECON-042}

{\sffamily\small\color{muted}\textbf{ECON-042} | R | Source-backed agenda; exact formalization is editorial\par}

\subsection*{Definitions and formal target}

This is a parameterized model-building and empirical research question, not a specified numerical instance. On a common measurable history space with observable filtration \(\mathcal I_t\), let \(\mathcal P\ne\varnothing\) describe the candidate uncertainty laws for climate and technology. An instance must specify whether these are laws of primitive innovations or controlled laws of endogenous histories. In either representation, specify how each policy \(\pi\) and candidate \(P\) induce a probability law \(Q^{\pi,P}\) of the full history, including climate feedback and technology choices. The expectation below is under this induced law; this notation does not require a particular transition model.

Specify initial conditions, policy action spaces and timing, observations, measurable transition rules, economic primitives and spatial equilibrium conditions. A policy is nonanticipative with respect to \(\mathcal I_t\). Specify an equilibrium selection across policies and candidate laws, with the measurability needed for the displayed expectations. Define the common feasible class \(\Pi\) using resource, budget and physical constraints, stating whether they hold on every admissible history or almost surely under each \(Q^{\pi,P}\). In the latter case state how law-dependent null sets are treated. Report empty feasibility or failure of the required equilibrium existence separately from an optimization value; no policy recommendation is then defined.

For a finite horizon \(T\in\{0,1,2,\ldots\}\), discount \(\beta\in(0,1]\), and a declared fixed-person or fixed-location welfare account, let \(W_t^{\pi,P}\) be the measurable real-valued aggregate welfare at date \(t\) in the selected equilibrium. The regional expression \(W_t^{\pi,P}=\sum_r\omega_rN_{rt}^{\pi,P}u(c_{rt}^{\pi,P})\) is valid only when it represents that account: \(r\) ranges over the specified finite region set, \(N_{rt}^{\pi,P}\geq0\) is population mass, \(c_{rt}^{\pi,P}\) lies in a declared consumption domain \(C\), \(u:C\to\R\) has a declared normalization, and \(\omega_r\) are declared welfare weights. Specify consumption units, price normalization and weight signs and normalization. A zero-population region contributes zero without requiring an average consumption there. Utility of average consumption must not replace aggregate individual utility without an aggregation justification. For heterogeneous people, define welfare by the corresponding sum or population integral of individual utilities and weights; fixed-person weights follow identities through migration, whereas fixed-location weights follow residence. Declare population membership and the treatment of migration, births, deaths and any welfare outside the living population. These choices are inputs, not conventions selected by [S1]. For each such specified instance with nonempty \(\Pi\), consider the editorial robust objective
\[
 \sup_{\pi\in\Pi}\inf_{P\in\mathcal P}
 \mathbb E_{Q^{\pi,P}}\sum_{t=0}^{T}\beta^t
                W_t^{\pi,P}.
\]
Specify and empirically discipline \(\mathcal P\) using declared observations, their measurement and sampling model, and stated structural restrictions. Explain how these inputs construct or assess the candidate set, distinguishing calibration, identification, sampling uncertainty and ambiguity; no particular statistical threshold is prescribed here. Declare a single-law benchmark \(P_0\), whether it belongs to \(\mathcal P\), and its induced laws. Specify which policy instruments count as mitigation or adaptation and which comparison outputs are sought: policy paths, welfare values, or cross-evaluation of policies under declared laws. State which feasible sets, equilibrium selections and welfare accounts are shared and identify any differences. Completion of an instance requires the specified model and empirical justification together with the declared comparisons; merely writing a robust objective does not supply them.

\subsection*{Clarifications and required assumptions}

The infimum is over alternative laws, not a draw of \(P\) from a distribution. Distinguish probabilistic risk from ambiguity about its law. Define support, correlations, damage uncertainty, available observations, boundedness/integrability, equilibrium selection, and the meaning of feasibility under all laws. Declare whether policy and robust values are real or extended real, ensure each expectation exists without an undefined difference of infinities, and provide bounds sufficient for any asserted finite value. Pointwise integrability alone does not ensure a finite robust value. Claims about arbitrarily near-optimal policies with an additive positive tolerance require a finite supremum value and nonempty feasibility. Ex-ante commitment max--min need not be dynamically consistent; recursive worst-case updating requires additional conditions such as a stated rectangular ambiguity structure. Optimizers require attainment assumptions. An infinite horizon additionally needs convergence.

\subsection*{Variants and changes of scope}

\begin{itemize}
\item Bayesian expected welfare, risk-sensitive utility and worst-case ambiguity are different objectives.
\item Learning can update a prior or an ambiguity set, but the update rule must be given.
\item Fixed person weights differ from weights tied to current residence after migration.
\end{itemize}

\subsection*{Audit conclusion and source relationship}

[S1], Section 7.4, printed p. 299, calls for uncertainty in spatial climate analysis. Max--min control is an editorial variation, not a policy criterion endorsed or solved by the source.

\subsection*{References and exact source roles}

{\footnotesize\raggedright
\par\noindent\textbf{[S1]} Klaus Desmet and Esteban Rossi-Hansberg (2024). \emph{Climate Change Economics over Time and Space}. Annual Review of Economics 16, 271–304. DOI: 10.1146/annurev-economics-072123-044449.\par \textit{Verified locator / source role:} Section 7.4, Uncertainty, p. 299.\par \url{https://rossihansberg.economics.uchicago.edu/CCETS.pdf}\par\smallskip
}

\subsection*{Common conventions: Additional-source status and mathematical conventions}

\label{audited:conventions}

Of the 100 supplied ECON entries, 65 distinct problems appear here; 32 close
matches refine earlier entries, and three duplicate questions map to existing
entries. Appendix~\ref{app:audited-crosswalk} lists every destination.
The attachment's P/R/S/U distinctions and source audits are retained. Its
precise mathematical formalizations are editorial unless the individual audit
says otherwise. Candidate subsidy targets ECON-004--006 retain their U flags;
the resolved approval-core question ECON-007 updates OP-0202 above.
The following supplied conventions govern both this part and the attributed
ECON refinements elsewhere in the catalogue, subject to local restrictions.

\textbf{Parameterized questions.} A problem family lists inputs that must be fixed before an instance is evaluated: population, horizon, policies, information, data law, model class and welfare weights. These are required choices, not quantities the citations are assumed to specify. Undefined applications should not be treated as identified estimands.

\textbf{Indices and integrability.} Sets of agents, firms and regions are finite unless a population measure is explicitly used. \(i,j\) index units, \(r\) regions, \(t\) dates and \(h\) horizons. \(\mathbb E\) and \(\Pr\) refer to the declared population and law; \(\mathbf1\{A\}\) is an indicator. Every displayed expectation and discounted sum needs existence in the stated domain. Infinite horizons use a discount in \((0,1)\) and a convergence condition; finite horizons may use discount one.

\textbf{Optimization and existence.} A displayed \(\max\), \(\min\), or \(\arg\max\) is conditional on attainment. For nonempty finite feasible menus with defined objective values this is immediate; general spaces need continuity and compactness/coercivity or another existence argument. Without attainment, the target is the supremum/infimum and arbitrarily near-optimal feasible choices. \(\inf\varnothing=+\infty\). Empty feasibility and an unidentified target are different issues.

\textbf{Potential outcomes.} \(Y_i(z)\) is the outcome under a specified intervention, not the observed conditional mean among people choosing \(z\). In binary comparisons \(\Delta Y_i=Y_i(1)-Y_i(0)\). Specify assignment, treatment versions, compliance, information and observation horizon. Interference is allowed under a full policy vector; compressed exposure notation needs a justified mapping. No interference, consistency and overlap are substantive assumptions, not automatic consequences of notation.

\textbf{Populations and observation.} Fix baseline eligibility and population weights. Conditioning on post-policy employment, survival, relocation or program uptake can change the population being compared. Death is not ordinary loss to follow-up; outcomes truncated by death need a composite convention, joint survival target, or explicitly defined survivor stratum. Fertility-changing interventions need a population functional rather than an unexplained fixed descendant cohort. Measurement and missingness models must be supplied.

\textbf{Identification.} A structural model \(m\in\mathcal M\) implies observable law \(P_m\) and target \(\theta(m)\). Identification means
\[
 P_{m_1}=P_{m_2}\ \Longrightarrow\ \theta(m_1)=\theta(m_2)
 \quad\text{for all }m_1,m_2\in\mathcal M .
\]
Without identification, the target is
\[
 \Theta(P;\mathcal M)=\{\theta(m):m\in\mathcal M,\ P_m=P\}.
\]
Writing \(\theta(P)\) before establishing identification can hide incompatible latent models. Confidence sets additionally need a sampling law, confidence level, asymptotic sequence and uniformity class. Point identification, coverage and informative precision are separate properties.

\textbf{Mediation.} Total effects of randomized assignment are distinct from cross-world natural indirect effects. Belief/effort-channel effects require a meaningful mediator intervention and additional measurement, exclusion or mediation assumptions. Randomization of the initial policy alone does not supply those assumptions. Report bounds or interventional alternatives when the stronger target is unsupported.

\textbf{Equilibrium.} A counterfactual \(W(z)\), \(p(z)\), or \(\mu_t^z\) requires individual optimization, feasibility, government/resource budgets, market clearing and state-transition laws. State equilibrium existence and selection, or report ranges over compatible equilibria. Initial-state integration includes subsequent endogenous transitions; current-state integration must use the policy-dependent distribution. Reduced-form invariance after policy is not assumed.

\textbf{Welfare accounts.} Scores, utility, health, dollars and ecological quantities require explicit conversion before addition. Fees, taxes, subsidies, wages and extortion payments are transfers between parties; distributional weights can make them welfare relevant, but their face value is not automatically a resource loss. State ownership, geographic boundaries, foreign-income weights and fiscal finance. Do not add profits to household welfare already including dividends, or count land capitalization and the underlying benefit twice. A benefit-cost expression must distinguish gross benefits from net welfare and subtract every real cost exactly once.

\textbf{Derivatives and risk.} Log derivatives require positive arguments; local responses require admissible perturbations and specified equilibrium branches. Differentiation under expectations needs regularity/domination, and boundaries may allow only one-sided derivatives. Risk uses a specified law; ambiguity uses a declared nonempty law set on a common history space. Adaptive policies must be nonanticipative. A robust criterion is an editorial choice unless attributed explicitly.

\textbf{Scope overlaps.} Allocation, collective choice, contracts, household bargaining and coordinated policy overlap economics with optimization and game theory. They are retained because they appeared in the original catalogue; no claim of a strictly game-theory-free selection is made.
% END ECONOMICS STATEMENT
\end{document}
